\section{Discussion}
\label{sec:discussion}

\paragraph{Direct against iterated.} V2 is an iterated forecast: it
estimates daily dynamics and simulates them forward. B-EQ is a direct one:
it takes each horizon's change from the account's past windows of that
length. Iterated forecasts are the more efficient when the daily model is
correctly specified, and direct forecasts are the more robust when it is
not \citep{marcellino2006direct}. Our results fit the second half of that
statement: V2 has lower CRPS at \hA\ days, where errors in the daily model
have few steps to compound, and B-EQ wins at \hB\ and \hC\ days. On the
macroeconomic series of \citet{marcellino2006direct} the iterated method
usually won; on sparse personal transactions, in our data, it wins only at
the shortest horizon. We have not shown which part of V2's daily model is
misspecified; \S\ref{sec:explanation} rules out some candidates.

\paragraph{Dispersion growth (a hypothesis, untested).} From \hA\ to \hC\
days, V2's median \lC\% width grows by a factor of \R{d/wg/V2/90}
(\R{d/wg/V2/30} by \hB\ days), and B-EQ's by \R{d/wg/B-EQ/90}
(\R{d/wg/B-EQ/30}). A random walk's width grows with $\sqrt h$, a factor
of \R{d/wg/rw/90} (\R{d/wg/rw/30}); linear growth would give
\R{d/wg/lin/90}. These are ratios of medians over accounts, not per-account
ratios. In V2's sparse branch, each path draws a rate $\lambda$ and then
signed amounts $A$, so the $h$-day sum $S_h$ has
\[
\operatorname{Var}(S_h) = h\,\mathbb E[\lambda]\,\mathbb E[A^2]
  + h^2 \operatorname{Var}(\lambda)\,(\mathbb E[A])^2 ,
\]
whose second term grows quadratically with the horizon. With
$\lambda \sim \mathrm{Gamma}(n + \tfrac12, \text{rate} = D)$,
$\operatorname{Var}(\lambda)/\mathbb E[\lambda] = 1/D$, so the ratio of
the quadratic to the linear term is
\[
R_h = \frac{h}{D}\,\frac{(\mathbb E[A])^2}{\mathbb E[A^2]} \le \frac{h}{D}.
\]
For a full \R{c/sparse-window}-day window this bounds the extra variance
of the discretionary flows at \R{d/rh/var/7}\%, \R{d/rh/var/30}\% and
\R{d/rh/var/90}\% at \hA, \hB\ and \hC\ days (\R{d/rh/sd/7}\%,
\R{d/rh/sd/30}\% and \R{d/rh/sd/90}\% in standard deviation). The bound
holds only for the discretionary component under the model's own
assumptions, not for the full forecast's quantile widths, and a shorter
window gives a larger bound. Four candidate explanations, none tested:
\begin{enumerate}[label=(\alph*)]
\item \textbf{Shared rate uncertainty.} The rate drawn once per path adds
  the $h^2$ term; by the bound above, it is unlikely to explain most of
  the gap.
\item \textbf{Balance-dependent behaviour.} Real balances may spread more
  slowly than a random walk, for instance because spending responds to the
  balance; B-EQ's historical windows would inherit this, and V2's additive
  simulation cannot represent it.
\item \textbf{Paycheck-state diffusion.} V2's expected paycheck is a random
  walk, so in the simplified case of equally spaced paydays the variance
  of cumulative income grows roughly with the cube of the number of
  paydays.
\item \textbf{B-EQ too narrow at long horizons.} A
  \R{c/lookback}-day lookback holds few non-overlapping \hC-day windows,
  which may make B-EQ understate long-horizon dispersion, consistent with
  its weaker \hC-day coverage on accounts with under a year of history
  (\R{ho/str/history-lt-12-months/B-EQ/90/cov90}\% against V2's
  \R{ho/str/history-lt-12-months/V2/90/cov90}\%, \S\ref{sec:reversal}).
\end{enumerate}
The v2 plan tests each of them (\texttt{paper/REVISION.md}).
